% =====================================================================
%  Chương 4 — Kết luận và Hướng phát triển
% =====================================================================
\chapter{Kết luận và Hướng phát triển}
\label{ch:ketluan}

\section{Tổng kết kết quả}
\label{sec:tongket}

Hai hệ thống thông minh đã được xây dựng trọn vẹn, mỗi hệ thống đi hết chuỗi
sáu tầng ở Hình~\ref{fig:sodo-hethong}: từ dữ liệu thật, qua biểu diễn véc-tơ,
qua năm mô hình học máy được so sánh dưới cùng một giao thức, tới dự đoán, và
cuối cùng tới một khuyến nghị hành động \emph{có chi phí xác định}.

\begin{table}[H]
\centering
\caption{Tổng kết hai hệ thống}
\label{tab:tongket}
\small
\begin{tabularx}{\textwidth}{@{}L{3.4cm}XX@{}}
\toprule
 & \textbf{Hệ thống 1 --- Tiểu đường} & \textbf{Hệ thống 2 --- Bất động sản} \\
\midrule
Loại bài toán   & Phân lớp nhị phân & Hồi quy \\
Quy mô dữ liệu  & 768 hồ sơ & 30.229 tin đăng \\
Số chiều biểu diễn & 12 & 84 \\
Mô hình cơ sở   & Accuracy $0{,}6547$, Recall $0{,}0000$ & MAE $1{,}8587$, $R^2 = -0{,}0008$ \\
Mô hình tốt nhất & K-Nearest Neighbors, $F_1 = 0{,}6809$ & XGBoost, $R^2 = 0{,}6141$ \\
Độ đo phụ tốt nhất & ROC-AUC $0{,}8709$ (Logistic Regression) & MAE $1{,}0493$ tỷ VNĐ \\
Cải thiện so với cơ sở & $F_1$: $0 \to 0{,}68$ & MAE giảm $43{,}5\%$ \\
Thuật toán hay biểu diễn chi phối & \textbf{Thuật toán}, gấp 6 lần & \textbf{Ngang bằng} \\
Đồ thị tri thức & 5 tầng, gói chăm sóc có chi phí & 5 tầng, gói vay có lịch trả nợ \\
\bottomrule
\end{tabularx}
\end{table}

Về hạ tầng: một lệnh \texttt{python run\_pipeline.py} tái tạo toàn bộ trong
khoảng 90 giây; 355 bài kiểm thử tự động chạy trong 17 giây; Đồ thị Tri thức
gồm 117 thực thể, 133 cạnh và 28 loại quan hệ đã nạp thật lên máy chủ Neo4j
Aura~\cite{neo4j}; ứng dụng web~\cite{streamlit} đã triển khai công khai và truy cập được từ internet~\cite{repo}.

\section{Phân tích khoa học --- bảy phát hiện}
\label{sec:phantich}

Mục này đáp ứng yêu cầu R10. Điểm chung của bảy phát hiện dưới đây: chúng đều
\emph{bác bỏ} một giả định mà người viết tin là đúng lúc mới bắt đầu, và mỗi
phát hiện đều kèm bằng chứng định lượng tái lập được.

\subsection{Phát hiện 1 --- Kết quả vượt trần tài liệu là dấu hiệu của lỗi}

Bản triển khai đầu tiên điền khuyết bằng trung vị theo nhóm \texttt{Outcome} và
đạt Accuracy $0{,}8705$ --- vượt hẳn trần $0{,}77$--$0{,}79$ mà mọi nghiên cứu
công bố trên tập Pima đều dừng lại. Nguyên nhân là rò rỉ nhãn: với
\texttt{Insulin} khuyết $374/768$ và \texttt{SkinThickness} khuyết $227/768$,
gần một nửa dữ liệu sau khi điền đã mang dấu vết của chính nhãn cần dự đoán.
Đưa \texttt{SimpleImputer} vào trong \texttt{Pipeline} đưa con số về $0{,}7914$.

\textbf{Bài học phương pháp:} đối chiếu kết quả với trần đã công bố trong tài
liệu là một bước kiểm tra rẻ tiền và hiệu quả bất ngờ. Một con số đẹp hơn dự
kiến đáng để đi tìm lỗi, không phải để ăn mừng.

\subsection{Phát hiện 2 --- Quét ngưỡng đòi hỏi mô hình cho xác suất liên tục}

Thực nghiệm quét ngưỡng chạy lần đầu trên Decision Tree sâu 6 tầng cho bảy
ngưỡng từ $0{,}50$ xuống $0{,}20$ ra \emph{đúng một} kết quả. Cây nông chỉ trả
về vài giá trị xác suất rời rạc nên phần lớn ngưỡng rơi vào cùng một khoảng.
Hiện tượng này để lại một dấu vết nhìn thấy được: đường ROC của Decision Tree ở
Hình~\ref{fig:dia-roc} có dạng gãy khúc thành vài đoạn thẳng dài thay vì cong
mượt như bốn mô hình còn lại. Chuyển sang XGBoost cho 12 mức kết quả khác nhau
trên 13 ngưỡng.

\textbf{Bài học:} một thực nghiệm không cho kết quả có thể là do thực nghiệm sai
thiết kế, không phải do hiệu ứng không tồn tại.

\subsection{Phát hiện 3 --- Cực trị chạm biên dải khảo sát thì chưa phải cực trị}

Bẫy này xuất hiện \textbf{hai lần độc lập}: độ sâu rừng ngẫu nhiên đạt đỉnh ở
$24$ khi dải quét dừng ở $24$, và ngưỡng quyết định đạt đỉnh ở $0{,}20$ khi dải
quét dừng ở $0{,}20$. Cả hai lần đều phải nới dải mới xác nhận được đỉnh thật
(lần lượt là $24$ và $0{,}15$).

\textbf{Bài học:} trước khi kết luận một siêu tham số là tối ưu, hãy kiểm tra nó
không nằm ở đầu mút dải khảo sát. Kiểm tra này rẻ và tự động hoá được --- hàm vẽ
biểu đồ quét ngưỡng trong dự án đã được sửa để tự in cảnh báo ``chạm biên dải
quét'' khi điều đó xảy ra.

\subsection{Phát hiện 4 --- ``Biểu diễn quan trọng hơn thuật toán'' là SAI nếu phát biểu trần}
\label{sec:phanchieu}

Đây là phát hiện trung tâm của cả báo cáo, và nó đáp ứng trực tiếp yêu cầu R13.

Cùng một thiết kế Thực nghiệm~3 --- giữ nguyên thuật toán, chỉ đổi biểu diễn ---
chạy trên hai bài toán cho hai kết luận trái ngược:

\begin{table}[H]
\centering
\caption{Biên độ do biểu diễn so với biên độ do thuật toán, hai bài toán}
\label{tab:biendo-ca-hai}
\small
\begin{tabularx}{\textwidth}{@{}L{5.4cm}R{3.0cm}R{3.2cm}@{}}
\toprule
 & \textbf{Tiểu đường ($F_1$)} & \textbf{Bất động sản ($R^2$)} \\
\midrule
Biên độ do đổi \textbf{biểu diễn}  & 0{,}0300 & 0{,}2563 \\
Biên độ do đổi \textbf{thuật toán} & 0{,}1809 & 0{,}2584 \\
\midrule
Bên nào chi phối & \textbf{Thuật toán} (gấp 6) & \textbf{Ngang bằng} \\
\bottomrule
\end{tabularx}
\end{table}

\textbf{Điều phân biệt hai trường hợp không nằm ở thuật toán, mà nằm ở chỗ phép
biểu diễn có làm thay đổi \emph{lượng} thông tin hay chỉ đổi \emph{dạng} thông
tin đã có.}

\begin{itemize}
  \item Ở bài toán bất động sản, phương án ``chỉ đặc trưng định lượng''
        \textbf{bỏ đi} toàn bộ nhóm biến định danh --- và \texttt{Province} cùng
        \texttt{District} không suy ra được từ bất kỳ cột nào khác. Đây là
        \emph{thêm hoặc bớt thông tin}, nên tác động lớn: $R^2$ mất $0{,}2563$.
  \item Ở bài toán tiểu đường, cả bốn phương án đều dùng đúng tám cột gốc ấy;
        chúng chỉ khác nhau ở chỗ có chuẩn hoá hay không, có nhân hai cột với
        nhau hay không. Đây là \emph{đổi dạng} thông tin đã có, nên tác động nhỏ
        và nằm trọn trong nhiễu lấy mẫu: bốn phương án chỉ phân loại khác nhau
        \textbf{1--9 ca trên 139}.
\end{itemize}

\begin{phathien}{Phát biểu đã được sửa lại cho đúng}
Phát biểu ban đầu --- ``biểu diễn dữ liệu quan trọng hơn thuật toán'' --- không
đứng vững qua kiểm chứng. Phát biểu đúng là:

\begin{quote}
\itshape
Biểu diễn dữ liệu chi phối kết quả khi và chỉ khi nó \textbf{thay đổi lượng
thông tin} có trong biểu diễn. Khi nó chỉ đổi \textbf{dạng} của thông tin đã có
--- chuẩn hoá, rời rạc hoá, nhân hai cột --- thì tác động của nó thường nhỏ hơn
tác động của việc đổi thuật toán, và có thể nằm trọn trong nhiễu lấy mẫu.
\end{quote}

Đây là lý do phải chạy \emph{cùng một thực nghiệm trên hai miền dữ liệu}: một
miền đơn lẻ sẽ cho một kết luận nghe rất thuyết phục và không có cách nào phát
hiện ra nó chỉ đúng một nửa.
\end{phathien}

\subsection{Phát hiện 5 --- ``Đặc trưng tốt'' là thuộc tính của cặp đặc trưng -- thuật toán}

Đặc trưng \texttt{Glucose\_BMI\_Risk} đồng thời là:

\begin{itemize}
  \item đặc trưng \textbf{tương quan mạnh nhất} với nhãn ($\rho = 0{,}5199$,
        vượt cả \texttt{Glucose} gốc $0{,}4928$);
  \item đặc trưng \textbf{quan trọng nhất} của XGBoost ($0{,}1874$);
  \item và là đặc trưng \textbf{làm Logistic Regression kém đi}
        ($F_1$ giảm từ $0{,}6818$ xuống $0{,}6667$).
\end{itemize}

Nguyên nhân của điều thứ ba là đa cộng tuyến dưới phạt $L_2$: đặc trưng này
tương quan $0{,}82$ với \texttt{Glucose} và $0{,}73$ với \texttt{BMI}, nên số
hạng $\lambda\lVert\bm\theta\rVert_2^2$ ở~\eqref{eq:regularized} phân tán trọng
số ra ba cột gần trùng nhau.

Hiện tượng đối xứng xảy ra ở phía kia: \texttt{Glucose\_Level} --- bản rời rạc
hoá của \texttt{Glucose} --- xếp \emph{cuối bảng} mức quan trọng của XGBoost
($0{,}0330$), vì cây quyết định tự tìm được ngưỡng chia tối ưu và việc rời rạc
hoá thủ công chỉ giới hạn không gian tìm kiếm của nó. Cùng một phép biến đổi có
ích cho mô hình tuyến tính và có hại cho mô hình cây.

\textbf{Bài học:} không tồn tại ``bộ đặc trưng tốt nhất'' độc lập với thuật
toán. Kỹ nghệ đặc trưng và chọn mô hình là một bài toán \emph{chung}, không phải
hai bước nối tiếp.

\subsection{Phát hiện 6 --- Giả định về phân phối phải được đo trước khi viện dẫn}

Giả định ``giá nhà lệch phải mạnh'' nghe hiển nhiên tới mức suýt trôi qua mà
không ai kiểm. Đo thật thì: \texttt{Price} có độ lệch $-0{,}029$, độ nhọn
$-0{,}8464$, trung bình $5{,}8721$ xấp xỉ trung vị $5{,}90$, và toàn bộ dữ liệu
nằm gọn trong dải $1{,}0$--$11{,}5$ tỷ. Tập dữ liệu đã bị cắt biên từ nguồn.

Giả định đúng với thị trường thật và \emph{sai} với tập dữ liệu này. Nếu tin
theo nó, ta sẽ áp phép biến đổi log cho biến mục tiêu --- và tạo ra độ lệch trái
thay vì sửa chữa gì.

Cùng bài học ấy lặp lại ở một chỗ khác: \texttt{Area} thì \emph{thật sự} lệch
phải mạnh (độ lệch $3{,}889$), và biến \emph{dẫn xuất} đơn giá mỗi m$^2$ cũng
vậy (độ lệch $2{,}135$, trung vị $102{,}3$ triệu đồng/m$^2$). Hình dáng mà trực
giác mong đợi có tồn tại trong dữ liệu --- chỉ là ở một biến khác với biến ta
tưởng.

\subsection{Phát hiện 7 --- Sai số gộp che mất hai chế độ sai khác hẳn nhau}

Mô hình định giá có MAE gộp $1{,}0493$ tỷ VNĐ. Tách theo ngũ phân vị giá thật
thì con số ấy vỡ ra thành hai chế độ:

\begin{itemize}
  \item Nhóm rẻ nhất: định giá \textbf{vượt} trung bình $1{,}0488$ tỷ, MAE
        $1{,}1575$.
  \item Nhóm giữa: gần như không chệch ($+0{,}2419$), MAE $0{,}7889$.
  \item Nhóm đắt nhất: định giá \textbf{hụt} trung bình $1{,}4844$ tỷ, MAE
        $1{,}5335$ --- gần gấp đôi nhóm giữa.
\end{itemize}

Độ chệch giảm \emph{đơn điệu} qua năm nhóm --- đây là quy luật, không phải nhiễu.
Nguyên nhân là hồi quy về trung bình cộng với việc mô hình cây không ngoại suy
ra ngoài khoảng giá trị đã thấy lúc huấn luyện.

\textbf{Hệ quả bắt buộc với người dùng:} với bất động sản cao cấp, con số hệ
thống đưa ra là \textbf{cận dưới}. Một hệ thống định giá im lặng về điều này sẽ
khiến người bán nhà cao cấp bị thiệt một cách có hệ thống --- và điều tệ nhất là
họ không có cách nào biết.

\section{Phản chiếu về vai trò của biểu diễn dữ liệu}
\label{sec:phanchieu-tong}

Yêu cầu R13 đòi hỏi một phản chiếu về biểu diễn dữ liệu. Ba tầng phản chiếu, đi
từ cụ thể tới khái quát:

\subsection{Tầng 1 --- Biểu diễn quyết định thuật toán nhìn thấy gì}

Cùng một hiện tượng thực tế cho những biểu diễn khác nhau, và mỗi biểu diễn mở
ra hoặc đóng lại một họ thuật toán:

\begin{itemize}
  \item Không chuẩn hoá thì K-Nearest Neighbors vô dụng, vì khoảng cách
        Euclid~\eqref{eq:euclid} bị cột \texttt{Insulin} thang trăm lấn át hoàn
        toàn. Nhưng chuẩn hoá lại \emph{không đổi được một ca nào} cho Logistic
        Regression --- hai phương án ``có'' và ``không có'' cho $F_1$ trùng khít
        tới bốn chữ số thập phân.
  \item Không mã hoá biến định danh thì $85{,}4\%$ năng lực quyết định của mô
        hình định giá biến mất, dù mọi cột định lượng vẫn còn nguyên.
\end{itemize}

\subsection{Tầng 2 --- Biểu diễn có thể mã hoá lẫn cả câu trả lời}

Đây là mặt tối của tầng 1. Bước điền khuyết ở Phát hiện~1 là một phép biến đổi
biểu diễn hoàn toàn bình thường về hình thức --- nó chỉ thay giá trị khuyết bằng
một con số hợp lý. Nhưng con số ấy được tính từ nhãn, nên biểu diễn thu được đã
chứa sẵn câu trả lời.

Điều nguy hiểm là hiện tượng này \emph{không} biểu hiện thành lỗi. Mã chạy trơn
tru, mọi bài kiểm thử xanh, và độ đo thì \emph{đẹp hơn} bình thường. Dấu hiệu
duy nhất là con số quá đẹp --- và nó chỉ thành dấu hiệu nếu ta biết trần tham
chiếu ở đâu.

\subsection{Tầng 3 --- Giới hạn của việc thiết kế biểu diễn}

Kết quả Thực nghiệm~3 ở cả hai bài toán cho thấy: khi thông tin đã có sẵn trong
dữ liệu, việc sắp xếp lại nó cho đẹp hơn mang lại rất ít. Bốn phương án biểu
diễn ở bài toán tiểu đường chỉ chênh nhau 1--9 ca trên 139; ba đặc trưng thiết
kế ở bài toán bất động sản làm $R^2$ giảm nhẹ chứ không tăng.

Cái mang lại nhiều là việc \textbf{đưa thêm thông tin mới vào}: biến định danh
về vị trí là thứ không cột nào khác suy ra được, và chính nó tạo ra biên độ
$0{,}2563$ $R^2$.

\begin{phathien}{Kết luận cuối cùng về biểu diễn dữ liệu}
Câu hỏi đúng khi thiết kế biểu diễn không phải \emph{``biểu diễn này có tinh vi
không''} mà là:

\begin{enumerate}
  \item \textbf{Nó có mang thêm thông tin mà dữ liệu gốc chưa có không?} Nếu
        không, hãy hạ kỳ vọng về mức cải thiện.
  \item \textbf{Nó có vô tình mang thông tin từ nhãn vào không?} Nếu có, mọi độ
        đo sau đó đều vô giá trị.
  \item \textbf{Nó hợp với thuật toán nào?} Chuẩn hoá bắt buộc với mô hình dựa
        trên khoảng cách, vô ích với mô hình tuyến tính, vô nghĩa với mô hình
        cây. Rời rạc hoá có ích cho mô hình tuyến tính, có hại cho mô hình cây.
\end{enumerate}
\end{phathien}

\section{Vai trò của Đồ thị Tri thức --- điều mô hình không làm được}
\label{sec:vaitro-kg}

Hai hệ thống trong báo cáo này có thể dừng lại ở tầng \textsc{Prediction} và vẫn
đáp ứng đủ yêu cầu về học máy. Việc đi thêm hai tầng nữa cho phép trả lời những
câu hỏi mà không độ đo nào chạm tới được:

\begin{table}[H]
\centering
\caption{Câu hỏi mà mô hình trả lời được và câu hỏi cần tới đồ thị tri thức}
\label{tab:vaitro-kg}
\small
\begin{tabularx}{\textwidth}{@{}L{4.0cm}XX@{}}
\toprule
 & \textbf{Mô hình trả lời} & \textbf{Đồ thị tri thức trả lời} \\
\midrule
Hệ thống Y sinh &
Xác suất mắc bệnh $54{,}5\%$ &
Mã ICD-10 E11; 5 mặt hàng cụ thể; chi phí ban đầu $1{,}99$--$2{,}75$ triệu đ;
duy trì $1{,}20$--$1{,}66$ triệu đ mỗi tháng; mua ở đâu; dịch vụ miễn phí nào đi
kèm \\
\addlinespace
Hệ thống Bất động sản &
Giá dự báo $8{,}42$ tỷ VNĐ &
Hạn mức vay $5{,}89$ tỷ; trả góp $47{,}5$ triệu đ/tháng; tổng lãi $8{,}34$ tỷ;
thu nhập tối thiểu $118{,}7$ triệu đ/tháng; bốn tiện ích trong bán kính $2$ km \\
\bottomrule
\end{tabularx}
\end{table}

\nhanxet{Khác biệt về \emph{loại} chứ không phải về \emph{mức độ}. Con số
$54{,}5\%$ không thể trở thành ``mua máy đo Accu-Chek giá $790$ nghìn'' bằng bất
kỳ phép tính nào trên dữ liệu huấn luyện --- tri thức ấy đến từ bên ngoài và
phải được khai báo tường minh.

Một điểm thiết kế đáng nói: nút \emph{gói} ở tầng 4 của cả hai đồ thị không phải
trang trí. Nó biến một danh sách rời rạc thành một đơn vị \textbf{có chi phí
tính được}, nhờ đó câu hỏi ``theo lời khuyên này thì tốn bao nhiêu'' có câu trả
lời ngay trên đồ thị chứ không phải cộng tay ở tầng giao diện. Cùng vai trò ấy,
nút \emph{trả góp} ở đồ thị bất động sản biến một hạn mức tĩnh thành một nghĩa
vụ tài chính kéo dài 300 kỳ.}

\section{Hạn chế của hệ thống}
\label{sec:hanche}

Báo cáo liệt kê hạn chế một cách thẳng thắn, vì một hệ thống được mô tả quá tốt
so với thực tế còn nguy hiểm hơn một hệ thống yếu được mô tả đúng.

\subsection{Hạn chế về dữ liệu}

\begin{enumerate}
  \item \textbf{Phạm vi dân số của tập Pima rất hẹp.} Toàn bộ đối tượng là nữ,
        từ 21 tuổi trở lên, thuộc cộng đồng người Pima --- nhóm có tỷ lệ mắc đái
        tháo đường thuộc hàng cao nhất thế giới. Áp dụng thẳng cho nam giới hoặc
        cho dân số Việt Nam là ngoại suy không có cơ sở.
  \item \textbf{Tập bất động sản đã bị cắt biên.} Hệ thống chỉ áp dụng được
        trong dải $1$--$11{,}5$ tỷ VNĐ. Ngoài dải ấy, mô hình cây vẫn trả về một
        con số nhưng con số đó vô nghĩa (Mục~\ref{sec:price-cat-bien}).
  \item \textbf{Không có yếu tố thời gian.} Cả hai tập là ảnh chụp tĩnh. Giá bất
        động sản biến động theo chu kỳ thị trường, nhưng mô hình không có cách
        nào biết tin đăng thuộc thời điểm nào.
  \item \textbf{Địa chỉ chỉ ở mức hành chính.} Không có toạ độ, nên không thể
        tính khoảng cách thật tới tiện ích. Đây là lý do cụm tiện ích ở
        Mục~\ref{sec:kg-hou} là mô hình hoá theo phân khúc chứ không phải khảo
        sát thực địa.
\end{enumerate}

\subsection{Hạn chế về mô hình}

\begin{enumerate}
  \item \textbf{Độ chệch đơn điệu theo phân khúc giá} --- Phát hiện~7. Hệ thống
        hiển thị cảnh báo nhưng không sửa được độ chệch này.
  \item \textbf{Trần hiệu năng của tập Pima.} Accuracy dừng ở mức
        $0{,}77$--$0{,}79$ trong toàn bộ tài liệu công bố, và tín hiệu phân tán
        mỏng trên nhiều đặc trưng (không đặc trưng nào vượt mức quan trọng
        $0{,}19$). Đây là giới hạn của dữ liệu, không phải của thuật toán.
  \item \textbf{Tập ngoại kiểm của bài toán tiểu đường quá nhỏ.} 77 ca không đủ
        để phân biệt hai mô hình chênh nhau dưới $0{,}05$ $F_1$
        (Mục~\ref{sec:tn4-dia}).
  \item \textbf{Không có ước lượng khoảng tin cậy.} Hệ thống trả về một con số
        điểm; một khoảng dự báo sẽ hữu ích hơn nhiều, đặc biệt ở phân khúc cao
        cấp nơi sai số lớn gấp đôi.
\end{enumerate}

\subsection{Hạn chế về tri thức miền trong đồ thị}

Ba ranh giới đã nêu ở Mục~\ref{sec:ranh-gioi-so-lieu} và được nhắc lại ở đây cho
đầy đủ: giá dược phẩm là khoảng tham khảo có ngày chốt chứ không phải bản ghi
giá trực tiếp; cụm tiện ích là mô hình hoá theo phân khúc; điều kiện vay là minh
hoạ theo mặt bằng lãi suất công bố chứ không phải chào giá của một ngân hàng cụ
thể. Riêng \emph{phép tính} trả góp, tổng lãi và thu nhập tối thiểu thì là toán
học thuần tuý, kiểm chứng được bằng công thức~\eqref{eq:nienkim} và~\eqref{eq:dti}.

\section{Hướng phát triển}
\label{sec:huongphattrien}

\begin{table}[H]
\centering
\caption{Hướng phát triển, sắp theo tỷ lệ giá trị trên công sức}
\label{tab:huongphattrien}
\small
\begin{tabularx}{\textwidth}{@{}L{4.2cm}XL{1.9cm}@{}}
\toprule
\textbf{Hướng} & \textbf{Nội dung và lý do} & \textbf{Công sức} \\
\midrule
Khoảng dự báo thay cho số điểm &
Dùng hồi quy phân vị hoặc conformal prediction để trả về khoảng
$[\hat y_{\text{thấp}}, \hat y_{\text{cao}}]$. Giải quyết trực tiếp hạn chế lớn
nhất hiện nay: người dùng không biết con số đáng tin tới đâu. &
Thấp \\
\addlinespace
Hiệu chỉnh độ chệch theo phân khúc &
Học một hàm hiệu chỉnh trên phần dư theo ngũ phân vị
(Bảng~\ref{tab:chech-hou}) để bù lại độ chệch đơn điệu. &
Thấp \\
\addlinespace
Bổ sung toạ độ địa lý &
Mã hoá địa chỉ thành kinh độ -- vĩ độ, tính khoảng cách thật tới tiện ích. Biến
cụm tiện ích từ mô hình hoá thành đo đạc. &
Trung bình \\
\addlinespace
Diễn giải mức cá thể (SHAP) &
Mức quan trọng hiện tại là mức toàn cục. SHAP cho biết \emph{với căn nhà này},
đặc trưng nào đẩy giá lên hay xuống bao nhiêu. &
Trung bình \\
\addlinespace
Yếu tố thời gian &
Thu thập dữ liệu có mốc thời gian để mô hình hoá chu kỳ thị trường. &
Cao \\
\addlinespace
Mở rộng đồ thị tri thức &
Nối thêm dữ liệu quy hoạch, hạ tầng giao thông đang xây, và lịch sử giao dịch
theo khu vực. &
Cao \\
\bottomrule
\end{tabularx}
\end{table}
